\section{Translation in terms of the fundamental group}
\label{IX.5}

Let
$$
g\colon S'\to S
$$
be an \emph{effective descent morphism} for the fibered category of
\emph{étale coverings} of preschemes, for example a proper, surjective
morphism of finite presentation \eqref{IX.4.12}, or a faithfully flat
and quasi-compact morphism. Introducing as usual~$S''$,~$S'''$, and
denoting by~$\cal{C}$,~$\cal{C}'$,~$\cal{C}''$,~$\cal{C}'''$ the
category of étale coverings of~$S$,~$S'$,~$S''$,~$S'''$ respectively,
we therefore have a $2$-exact diagram of categories
\begin{equation*}
\label{eq:IX.5.*}
\tag{$*$}
\xymatrix@C=.5cm{
{\cal{C}} \ar[r]^-{~p^*} & {\cal{C}'}
\ar@<2pt>[rr]^{p_1^*,p_2^*} \ar@<-2pt>[rr] && {\cal{C}''}
\ar@<4pt>[rrr]^{p_{21}^*,p_{32}^*,p_{31}^*} \ar[rrr] \ar@<-4pt>[rrr] &&&
{\cal{C}'''}
}
\end{equation*}
corresponding to the diagram
$$
\xymatrix@C=.5cm{ S & \ar[l]_-{~p} S' && \ar@<2pt>[ll] \ar@<-2pt>[ll]_-{p_1,p_2}
S'' &&& \ar@<4pt>[lll] \ar[lll] \ar@<-4pt>[lll]_{p_{21},p_{32},p_{31}} S'''.
}
$$
Suppose
% original p. 243
the preschemes~$S$,~$S'$,~$S''$,~$S'''$ are disjoint sums of
connected preschemes, which will be the case in particular if they
are locally connected preschemes, a fortiori if they are locally
noetherian (for example if~$S'$ is of finite type over~$S$ locally
noetherian). Then the categories~$\cal{C}$,~$\cal{C}'$ \ldots\ in~\eqref{eq:IX.5.*}
are multi-Galois categories V~\Ref{V.9}, each described by a
collection of totally disconnected compact topological groups,
namely the fundamental groups of the connected components of the
preschemes~$S$,~$S'$,~$S''$,~$S'''$. We assume for simplicity
that~$S$ is connected, and shall then give a procedure for computing
its fundamental group, in terms of the fibered category formed
with~$\cal{C}'$,~$\cal{C}''$,~$\cal{C}'''$, suitably made explicit by
means of the fundamental groups expressing these categories. The
reader will note that the procedure sketched is in fact valid in the
general framework of multi-Galois categories (which need not come
from given preschemes~$S$,~$S'$,~$S''$,~$S'''$). It is, moreover, the
analogue of the well-known procedure for computing the fundamental
group of a topological space~$S$ that is a locally finite union of
closed subspaces~$S_i$ (or an arbitrary union of open
subspaces~$S_i$), by means of the fundamental groups of the components
of the~$S_i$
and of the components of the $S_i\cap S_j$. Of course, the analogous
situation in the framework of preschemes does fall within the general
framework of descent, by introducing the prescheme~$S'$, sum of
the~$S_i$, and the canonical morphism $g\colon S'\to S$.

Set
$$
E'=\pi_0(S'), \quad E''=\pi_0(S''), \quad E'''=\pi_0(S'''),
$$
where~$\pi_0$ denotes the functor ``set of connected components''.
Since the fibered products of~$S'$ over~$S$ form a simplicial object
of~$\Sch$, it is transformed by the functor~$\pi_0$ into a simplicial
set of which~$E'$,~$E''$,~$E'''$ are the components of
dimension~$0$,~$1$,~$2$. We shall have to use the simplicial maps
$$
q_i=\pi_0(p_i),\quad(i=1,2)\quad\text{and}
\quad q_{ij}=\pi_0(p_{ij}),\quad (i,j)=(2,1),(3,2),(3,1),
$$
exhibited in the diagram
\begin{equation*}
\label{eq:IX.5.1}
\tag{1} {\xymatrix@C=.5cm{ E' && \ar@<2pt>[ll] \ar@<-2pt>[ll]_{q_1,q_2} E'' &&&
\ar@<4pt>[lll] \ar[lll] \ar@<-4pt>[lll]_{q_{21},q_{32},q_{31}} E'''. }}
\end{equation*}
The objects
% original p. 244
of~$E'$ will be denoted with a prime, like~$s'$, those of~$E''$
\resp $E'''$ will be denoted with a~$''$ \resp a~$'''$. The fact
that~$S$ is connected translates into $\pi_0(K)=0$, where~$K$ is the
simplicial set defined by $g\colon S'\to S$, or again into the fact
that the equivalence relation in~$E'$ generated by the pair of maps
$(q_1,q_2)$ is transitive.

We shall choose once and for all an element~$s_0'$ in~$E'$, and for
every~$s'$ in~$E'$, an element~$\overline{s'}\in E''$ such
that\footnote{Note that the element~$\overline{s'}$, whose existence
is implicitly assumed, does not exist in all cases. Consequently,
theorem~\Ref{IX.5.1} as it is stated does not apply in all cases. It
is not difficult, however, taking the written text as a guide, to
modify the statement of this theorem in such a way that it gives a
method of computation which applies in all cases. In particular, the
corollaries of the said theorem are valid as they stand.}
$$
q_1(\overline{s'})=s_0', \quad q_2(\overline{s'})=s',
$$
thus exhibiting the connectedness of~$S$. For every~$s'\in E'$, let
us choose a geometric point~$\underline{s}'$ in the connected
component~$s'$ of~$S'$; this point will in fact come in through the
corresponding fiber functor~$F_{s'}'$ on the multi-Galois
category~$\cal{C}'$. The group of automorphisms of this functor, \ie
the fundamental group of~$S'$ at~$\underline{s}'$, will be denoted
by~$\pi_{s'}$. We likewise choose points~$\underline{s}''$
and~$\underline{s}'''$, hence functors~$F_{s''}''$
and~$F_{s'''}'''$, whence fundamental groups~$\pi_{s''}$
and~$\pi_{s'''}$. Thus
$$
\pi_{s'}=\pi_1(S',\underline{s}'), \quad
\pi_{s''}=\pi_1(S'',\underline{s}''), \quad
\pi_{s'''}=\pi_1(S''',\underline{s}''').
$$
For every $s''\in E''$, $p_1(\underline{s}'')$ lies in the same
connected component as $\underline{q_1(s'')}$, hence there exists an
isomorphism of functors
$F''_{s''}\circ p_1^*\isomto F'_{s'}$
(\ie a ``class of paths'' from~$p_1(\underline{s}'')$ to
$\underline{q_1(s'')}$). The same remark holds for~$q_2$, and
the~$q_{ij}$. Let us choose all these classes of paths:
$$
F''_{s''}\circ p_i^*\isomto F'_{q_i(s'')},\quad F_{s'''}'''\circ
p_{ij}^*\isomto F'''_{q_{ij}(s''')},
$$
(for
% original p. 245
$i=1,2$ and $(i,j)=(2,1),(3,2),(3,1)$). In particular there result
homomorphisms of groups:
\begin{equation*}
\label{eq:IX.5.2}
\tag{2} {q_i^{s''}\colon \pi_{s''}\to\pi_{q_i(s'')},\quad
q_{ij}^{s'''}\colon \pi_{s'''}\to\pi_{q_{ij}(s''')},}
\end{equation*}
(same values of~$i$ and $(i,j)$). Finally, let us recall that the
structure of cloven category with fibers~$\cal{C}'$,~$\cal{C''}$,~$\cal{C'''}$
also contains isomorphisms of functors:
$$
p_{21}^*p_1^*\isomto p_{31}^*p_1^*,\quad p_{21}^*p_2^*\isomto
p_{32}^*p_1^*,\quad p_{31}^*p_2^*\isomto p_{32}^*p_2^*,
$$
deduced from isomorphisms of the two sides respectively with the
$u_i^*$~($i=1,2,3$), where the~$u_i$ are the three projections
of~$S'''$ into~$S'$. When one makes these data explicit, one finds,
for every~$s'''$, a well-determined element
\begin{equation*}
\label{eq:IX.5.3}
\tag{3} {a_i^{s'''}\in\pi_{v_i(s''')},}
\end{equation*}
(where the~$v_i$,~$i=1,2,3$ are the three maps $E'''\to E'$
defined by $v_i=\pi_0(u_i)$), subject moreover to the conditions:
$$
q_1^{s_1''}q_{21}^{s'''}=\mathrm{int}(a_1^{s'''})
q_1^{s_2''}q_{31}^{s'''}\qquad (s_1''=q_{21}(s'''),\
s_2''=q_{31}(s''')),
$$
and the two analogous conditions, involving the~$a_2$ and~$a_3$. The
reader will note, moreover, that the
data~\eqref{eq:IX.5.1},~\eqref{eq:IX.5.2},~\eqref{eq:IX.5.3} make it
possible to reconstruct, up to an equivalence of fibered categories,
the fibered category under consideration with
fibers~$\cal{C}'$,~$\cal{C}''$,~$\cal{C}'''$. They must therefore
make it possible in principle to reconstruct~$\cal{C}$ up to
equivalence, hence its fundamental group up to isomorphism. We shall
in fact determine the fundamental group at the geometric point
$p(\underline{s}_0')$ of~$S$, \ie the group of automorphisms
of~$F'_{s_0'}\circ p^*$.

We note that giving an object~$X'$ of~$\cal{C}'$ is essentially
equivalent to giving finite sets~$X'_{s'}$ ($s'\in E'$) on which
the~$\pi_{s'}$ operate continuously. A
% original p. 246
gluing datum on such an object then amounts to giving, for every
$s''\in E''$, a bijection:
$$
\varphi_{s''}\colon X'_{q_1(s'')}\isomto X'_{q_2(s'')}
$$
compatible with the operations of~$\pi_{s''}$, which operates on
both sides by means of the homomorphisms $q_i^{s''}\colon
\pi_{s''}\to\pi_{q_i(s'')}$. Taking first the~$s''$ of the
form~$\overline{s'}$, we see that such a datum defines bijections
$$
\psi_{s'}\colon X'_{s_0'}=F_0'(X')\to X_{s'}'
$$
which make it possible to identify the~$X'_{s'}$ with the same set
$F_0'(X')=X'_{s_0'}$, on which all the groups~$\pi_{s'}$ will
henceforth operate. This being so, the bijections~$\varphi_{s''}$
will correspond to bijections
$$
g_{s''}\colon F_0'(X')\isomto F_0'(X'),
$$
subject on the one hand to the relations of commutation
with~$\pi_{s''}$:
\begin{equation*}
\label{eq:IX.5.a}
\tag*{a)} {g_{s''}q_1^{s''}(g'')=q_2^{s''}(g'')g_{s''}\qquad (s''\in E'',\ g''\in\pi_{s''}),}
\end{equation*}
on the other hand to the relations
\begin{equation*}
\label{eq:IX.5.b}
\tag*{b)} {g_{\overline{s'}}=g_{\overline{s_0'}}\qquad (s'\in E'),}
\end{equation*}
expressing the way in which we had identified the~$X_{s'}$ with one
another. When one makes explicit the condition for such a gluing
datum to be in fact a descent datum, one finds the relations:
\begin{equation*}
\label{eq:IX.5.c}
\tag*{c)} {a_3^{s'''}g_{q_{31}(s''')}a_1^{s'''}=
g_{q_{32}(s''')}a_2^{s'''}g_{q_{21}(s''')}\qquad (s'''\in E''').}
\end{equation*}

This gives us an equivalence between the category of objects
of~$\cal{C}'$ endowed with a
% original p. 247
descent datum, and the category of finite sets on which the
groups~$\pi_{s'}$ operate continuously, endowed moreover with
bijections~$g_{s''}$ satisfying the
relations~\Ref{eq:IX.5.a},~\Ref{eq:IX.5.b},~\Ref{eq:IX.5.c}. Let
then~$G$ be the group generated by the groups~$\pi_{s'}$ and the new
generators~$g_{s''}$, subject to the
relations~\Ref{eq:IX.5.a},~\Ref{eq:IX.5.b},~\Ref{eq:IX.5.c}, and
let~$\pi$ be the group that is the projective limit of the quotients
of~$G$ by the subgroups of finite index whose inverse images in the
groups~$\pi_{s'}$ are open subgroups. We also say that~$\pi$ is the
\emph{group of Galois type generated by the~$\pi_{s'}$ and
the~$g_{s''}$, subject to the
relations~\Ref{eq:IX.5.a},~\Ref{eq:IX.5.b},~\Ref{eq:IX.5.c}}. One
verifies at once that the category under consideration is also
equivalent to the category of finite sets on which the topological
group~$\pi$ operates continuously. This establishes the following
statement:
\begin{theorem}
\label{IX.5.1}
Let $g\colon S'\to S$ be a morphism of preschemes which is an
effective descent morphism for the fibered category of étale
coverings of preschemes (\cf \Ref{IX.4.9} and~\Ref{IX.4.12}).
Suppose~$S$ connected, and~$S'$, its fibered square~$S''$ and its
fibered cube~$S'''$, sums of connected preschemes (which is the case
for example if~$S'$ is of finite type over~$S$ locally noetherian and
connected). Let us choose as above: a geometric point in every
connected component of~$S'$,~$S''$,~$S'''$, certain classes of
paths, an $s_0'\in E'$, and for every $s'\in E'$ an
$s''\in E''$
whose two images in~$E'$ are~$s_0'$ and~$s'$.
($E'$,~$E''$,~$E'''$ denote respectively the set of connected
components of~$S'$,~$S''$,~$S'''$.) Then the fundamental group
of~$S$ at the geometric point image of~$s_0'$ is canonically
isomorphic to the group of Galois type generated by the
$\pi_{s'}=\pi_1(S',\underline{s}')$ ($s'\in E'$) and
generators~$g_{s''}$ ($s''\in E''$), subject to the
relations~\Ref{eq:IX.5.a},~\Ref{eq:IX.5.b},~\Ref{eq:IX.5.c} above
involving the elements of the groups
$\pi_{s''}=\pi_1(S'',\underline{s}'')$, and the
elements~$a_i^{s'''}$ ($i=1,2,3$, $s'''\in E'''$) introduced above.
\end{theorem}

\begin{corollary}
\label{IX.5.2}
Suppose that~$S'$ and~$S''$ have only a finite number of connected
components, and that the fundamental groups of the connected
components of~$S'$ are topologically finitely generated. Then the
fundamental group of~$S$ is topologically finitely generated.
\end{corollary}

Thus,
% original p. 248
we shall prove later that the fundamental group of a normal
projective scheme over an algebraically closed field is topologically
finitely generated. Using Chow's lemma and the normalization of
algebraic schemes, it will follow that the same result is true for
every proper scheme over an algebraically closed field.

\begin{corollary}
\label{IX.5.3}
Suppose that~$S'$,~$S''$,~$S'''$ have only a finite number of
connected components, that the fundamental groups of the connected
components of~$S'$ are topologically \emph{finitely presented}, and
the fundamental groups of the connected components of~$S''$
topologically \emph{finitely generated}. Then the fundamental group
of~$S$ is topologically \emph{finitely presented}.
\end{corollary}

We note that one can express~\Ref{IX.4.9} (restricted to étale
\emph{coverings}) by saying that \emph{a finite radicial surjective
morphism of noetherian preschemes induces an isomorphism of the
fundamental groups}; figuratively, one can therefore say that the
fundamental group is a \emph{topological invariant} for preschemes.
More generally, one can make explicit, by means of~\Ref{IX.5.1}, the
effect on the fundamental group of operations on preschemes having a
simple topological meaning, such as the ``pinching'' of the prescheme
along a finite set of points. One finds for example:

\begin{corollary}
\label{IX.5.4}
Let $g\colon S'\to S$ be a finite morphism of finite presentation,
$T$ a discrete subset of~$S$. For every $s\in S$, let $n(s)$ be the
``geometric number of points'' in the fiber $g^{-1}(s)$ (which can
also be expressed as the separable degree of~$g^{-1}(s)$
over~$\kres(s)$, the sum of the separable degrees of its residue
extensions). Suppose that for $s\in S-T$ we have $n(s)=1$. For every
$s\in T$, let $K_s$ be an algebraically closed extension of
$\kres(s)$, $I_s$ the set of geometric points of~$S'$ with values
in~$K_s$ (it is a set with $n(s)$ elements), $I_s'$ the complement of
a chosen point of~$I_s$, and finally $I'$ the set that is the union of
the $I_s'$. Suppose $S'$ connected. Then the fundamental group of~$S$
is isomorphic to the group of Galois type generated by the
fundamental group of~$S'$ and generators $g_i$ ($i\in I'$), subject to
no additional condition.
\end{corollary}

The
% original p. 249
details of the proof are left to the reader; the statement obtained
is merely the translation, into the language of group theory, of the
fact that we have an equivalence between the category $C$ of étale
coverings of~$S$ and the category of étale coverings $X'$ of~$S'$
\emph{endowed}, for every $s\in T$, with a transitive system of
bijections between the $n(s)$ fibers of~$X'$ at the points
of~$g^{-1}(s)$ with values in~$K_s$. (In this intrinsic form, of
course, it is no longer necessary to assume $S'$ connected.)
\begin{example}
\label{IX.5.5}
One easily proves that the rational curve $\PP^1_k$ over an
algebraically closed field $k$ is simply
connected\footnote{\Cf Exp\ptbl XI~\Ref{XI.1.1}.}. Hence the
fundamental group of a complete rational curve having exactly one
double point, with $n$ analytic branches, is the free group of Galois
type generated by $n-1$ generators. For example, in the case of an
ordinary double point, one finds the fundamental group $\hat\ZZ$, as
announced in (I~\Ref{I.11}~a)). On the other hand, the existence of a
cusp (which is a ``geometrically unibranch'' point) has no influence
on the fundamental group.
\end{example}

\begin{corollary}
\label{IX.5.6}
Let $g\colon S'\to S$ be a universally submersive morphism of
preschemes with geometrically connected fibers, $S$ being connected.
Then $S'$ is connected, and choosing a geometric point $s'$ in~$S'$
and denoting by $s$ its image in~$S$, the homomorphism
$$
\pi_1(S',s')\to \pi_1(S,s)
$$
is \emph{surjective}. If $g$ is an effective descent morphism for the
fibered category of étale coverings of preschemes (\cf
\emph{\Ref{IX.4.12}}), then, introducing the geometric point
$s^{\prime\prime}=\diag(s')$ of~$S^{\prime\prime}=S'\times_S S'$, and
the two homomorphisms
$$
p_{1*},p_{2*}\colon \pi_1(S^{\prime\prime},s^{\prime\prime}) \to
\pi_1(S',s')
$$
induced by the two projections, $\pi_1(S, s)$ is isomorphic to the
cokernel of this pair of morphisms in the category of groups of Galois
type, \ie to the quotient of~$\pi_1(S',s')$ by the closed invariant
subgroup generated by the
% original p. 250
elements of the form
$p_{1*}(g^{\prime\prime})p_{2*}(g^{\prime\prime})^{-1}$, with
$g^{\prime\prime}\in \pi_1(S^{\prime\prime},s^{\prime\prime})$.
\end{corollary}

Indeed, we know by~\Ref{IX.3.4} that the functor $X\mto X\times_SS'$
from étale coverings over~$S$ to étale coverings over~$S'$ is fully
faithful, which is equivalent to the fact that the homomorphism on
the fundamental groups is an epimorphism (V~\Ref{V.6.9}). The last
assertion is an immediate consequence of the
description~\Ref{IX.5.1}.

\begin{remark}
\label{IX.5.7}
It is not known at present whether the fundamental group of a proper
scheme over an algebraically closed field $k$ is topologically
finitely presented\footnote{This seems very improbable in the case of
smooth curves of genus $g\geq 2$, in characteristic $p>0$. When one
replaces $\pi_1$ by its largest quotient prime to $p$, on the other
hand, it seems that the well-known techniques make it possible to
give an affirmative answer, even without a properness hypothesis.
\Cf a work in preparation by J.P\ptbl Murre.}. Using~\Ref{IX.5.3}, a
well-known technique of hyperplane sections, and the desingularization
of normal surfaces, one is reduced to the case of a \emph{smooth
surface over} $k$. This makes it possible at least to show, by
transcendental means, that the answer is affirmative in
characteristic $0$ (and this without being obliged to assume the
triangulability of singular algebraic varieties). In characteristic
$p>0$, the main difficulty seems to lie in the case of curves, of
which one knows only that the fundamental group is a quotient of the
one that occurs in the classical case (\cf the following exposé), the
kernel by which one divides being, however, very poorly known.
\end{remark}

\begin{remark}
\label{IX.5.8}
One could make explicit other special cases than~\Ref{IX.5.4}
and~\Ref{IX.5.6} where~\Ref{IX.5.1} takes a particularly simple form.
An interesting case is the one where $S$ is the quotient of~$S'$ by a
finite group of automorphisms $\Gamma$. \emph{Then the category of
étale coverings of~$S'$ is equivalent to the category of étale
coverings $X'$ of~$S'$ on which the group $\Gamma$ operates compatibly
with its operations on~$S'$, in such a way that for every $s'\in S'$
and every $g\in \Gamma_{s'}$} (where $\Gamma_{s'}$ denotes \emph{the
inertia group} of~$s'$ in~$\Gamma$), \emph{$g$~operates trivially on
the fiber} $X_{s'}^\prime$. If $S'$ is connected, this statement can
be interpreted as follows. Let $C'_0$ be the category of étale
coverings of~$S'$ on which $\Gamma$ operates compatibly with its
operations on~$S'$ (but without necessarily satisfying the above
condition on the inertia groups of the points of~$S'$). One easily
sees that it is a Galois category (V~\Ref{V.5}), and that for
% original p. 251
every geometric point $a'$ of~$S'$, the fiber functor $X'\mto X_{a'}'$
on~$C'_0$ is a fundamental functor. Let $\pi_1(S',\Gamma;a')=G$ be
the group of automorphisms of this functor, endowed with its usual
topology. We then have a canonical exact sequence
$$
e\to \pi_1(S',a')\to G\to \Gamma\to e
$$
(special case of (V~\Ref{V.6.13}), where one takes for $S$ the trivial
covering $S'\times \Gamma$ of~$S'$ defined by $\Gamma$, where one
lets $\Gamma$ operate in the obvious way). Moreover, for every
geometric point $b'$ of~$S'$, we have an isomorphism
$\pi_1(S', \Gamma ; b' )\to G = \pi_1(S', \Gamma; a')$ defined up to
an inner automorphism coming from~$\pi_1(S',a')$, and since
$\Gamma_{b'}$ maps in an obvious way into the left-hand side, we
obtain a homomorphism
$$
u_{b'}\colon \Gamma_{b'}\to G \quoi,
$$
defined up to an inner automorphism (coming from~$\pi_1(S',a')$),
whose composite with the canonical homomorphism $G\to \Gamma$ is,
moreover, the canonical immersion $\Gamma_{b'} \to \Gamma$. This
being so, \emph{the fundamental group $\pi_1(S,a)$ \emph{is}
canonically isomorphic to the quotient group of~$G = \pi_1(S',
\Gamma; a')$ by the closed invariant subgroup generated by the images
of the homomorphisms $\Gamma_{b'} \to G$}. In particular, the image
of~$\pi_1(S',a')$ in $\pi_1(S,a)$ is an invariant subgroup, and the
corresponding quotient is isomorphic to a quotient of~$\Gamma$. One
can, moreover, reduce the number of ``relations'' introduced by
introducing, for every $g\in \Gamma$, $ g\neq e$, the
subprescheme $S'_g$ of coincidences of the automorphisms $\id_S$ and
$g$ of~$S$, by choosing a geometric point $b'_{g,i}$ in each connected
component of~$S'_g$, and then one of the corresponding homomorphisms
$\pi_1(S',\Gamma; b'_{g,i})\to G$, whence liftings $\bar{g}_i$ of~$g$
in~$G$. It then suffices to take the quotient of~$G$ by the closed
invariant subgroup of~$G$ generated by the $\bar{g}_i$.

When $a'$ is invariant under $\Gamma$, one easily sees that $\Gamma$
operates in a natural way on~$\pi_1(S',a')$, and $G$ is identified
with the corresponding semidirect product. Then identifying~$\Gamma$
with a subgroup of~$G$, one sees that in the relations
% original p. 252
introduced above, taking $b'=a'$, one finds ``$g = e$'' for $g\in I$.
Hence \emph{if $S'$ has a geometric point $a'$ fixed by}~$I$ (\ie a
point $s'$ whose inertia group is $\Gamma$), \emph{then $\pi_1(S,a)$
is a quotient group of the quotient group of Galois type
of~$\pi_1(S',a')$ obtained by \emph{``making trivial''} the operations
of~$\Gamma$ on~$\pi_1(S',a')$; and it is even isomorphic to this
latter group if one assumes that for every $g\in G$, the inertia set
$S_g'$ is connected, hence passes through the locality of~$a'$}. This
last assertion is indeed contained in the second description given
above for the relations to be introduced in $G$.

This last result applies in particular if one takes for $S'$ the
cartesian power $X^n$ of a connected prescheme over an algebraically
closed field, for $\Gamma$ the symmetric group $\Gamma={\goth {S}}_n$,
operating in the usual way, whence for $S$ the $n$th symmetric power
of~$S$. Taking then for $a'$ a geometric point localized at the
diagonal, we are under the preceding conditions, the inertia sets
$S_g'$ indeed all containing the diagonal. Using the fact, proved in
the following exposé, that if $X$ is proper and connected over~$k$,
the fundamental group of~$X^n$ is identified with $\pi_1(X)^n$, we
find the following amusing result: \emph{If $X$ is proper and
connected over~$k$ algebraically closed, the fundamental group of its
$n$th symmetric power, $n \geq 2$, is isomorphic to the fundamental
group of~$X$ made abelian}. (I do not know whether the analogous fact
in algebraic Topology is known; it should be possible to establish it
by the same method of descent.) Taking for example for $X$ a rational
curve $X=\PP^1_k$, we find an $N$th proof of the fact that $\PP^r_k$
is simply connected, using the fact that $\PP^1_k$ is. Let us now
take for $X$ a simple curve over~$k$, and $n\geq 2g-1$, so that
$\Symm^n(X)$ is fibered over the jacobian $J$, with projective spaces as
fibers, hence (as we shall see by means of the results of the two
following exposés) has the same fundamental group as $J$. We then
recover without dévissage the well-known fact that \emph{the
fundamental group of the jacobian of~$X$ is isomorphic to the
fundamental group of~$X$ made abelian}.
\end{remark}
